\documentclass[11pt]{article}
\usepackage[margin=1in]{geometry}
\usepackage{amsmath,amssymb,amsthm}
\usepackage{array,longtable,booktabs}
\usepackage[hidelinks]{hyperref}

\newtheorem{theorem}{Theorem}
\newtheorem{proposition}[theorem]{Proposition}
\newtheorem{corollary}[theorem]{Corollary}
\newtheorem{example}[theorem]{Example}
\newcommand{\Ann}{\operatorname{Ann}}
\newcommand{\diag}{\operatorname{diag}}
\newcommand{\GL}{\operatorname{GL}}

\title{P6-T05: Signature, Nondegeneracy, and Source Covariance Remain Unselected}
\author{The AEther-Flow Research Project}
\date{2026-07-26}

\begin{document}
\maketitle

\begin{center}
\fbox{\parbox{0.92\textwidth}{
\textbf{Status and authority.}
This is a task-local \emph{draft/control} result about proposal-only
source-extension burdens. It proves exact comparison countermodels, a scoped
source-naturality obstruction, and a conditional coframe-gluing theorem. It
is not a canonical-ontology edit, source-law adoption or rejection, physical
metric or causal-cone construction, Gate Chair verdict, global no-go theorem,
publication, or completed derivation.
}}
\end{center}

\section{Question and fixed input}

P6-T03 leaves as its complete source-side geometric input
\[
  \mathcal D_1(S,V)=\bigl([V]_+,\Ann(V),\preceq_V\bigr),
\]
and proves that this one-ray datum does not select a Lorentzian conformal
class. P6-T04 tests the enlarged package
\[
  \mathcal D_{\rm cal}
  =\bigl(\mathcal D_1,\mu,\{\theta_A\}_{A\in\mathcal A}\bigr),
\]
where \(\mu\) is an independent positive source density and the
\(\theta_A\) are dimensionless scalar phases. It proves that these carriers
do not select conformal shape, physical metric scale, absolute clocks, or
rods.

P6-T05 asks whether the same package nevertheless fixes three remaining
mathematical obligations:
\begin{enumerate}
\item a Lorentzian rather than Euclidean or split signature;
\item nondegeneracy separated from a singular boundary; and
\item a natural transformation law under source-induced frame changes.
\end{enumerate}
No target atlas, target metric, GR null cone, matter response, \(G\), \(c\),
operational readout normalization, or benchmark fit is admitted. The historical scoped
\(g_{\rm eff}\) source-extension object is not a premise and its status is
unchanged.

\section{Typed assumption audit}

\begin{longtable}{>{\raggedright\arraybackslash}p{0.10\textwidth}
                  >{\raggedright\arraybackslash}p{0.38\textwidth}
                  >{\raggedright\arraybackslash}p{0.41\textwidth}}
\toprule
ID & Test assumption & Status and consequence\\
\midrule
\endhead
A1 & \(S\) is smooth and four dimensional and \(V\) is smooth, nonzero, and
positively oriented. & Inherited candidate scope; not a physical spacetime or
clock conclusion.\\
A2 & Only \(\mathcal D_{\rm cal}\) is supplied unconditionally. & No hidden
coframe, bilinear form, connection, target chart, or metric-volume identity.\\
A3 & \(\mu\) is independent source density data. & It is not assumed equal to
the volume density of any comparison form.\\
A4 & Scalar phases may be chosen constant on the exact countermodel datum. &
They then add no frame discriminator, while remaining valid admitted data.\\
A5 & A reconstruction rule must be source-natural: automorphisms of its input
must act equivariantly on its output. & This is a consistency burden, not an
extra geometric premise.\\
A6 & Comparison forms are mathematical controls only. & Their signatures,
cones, and volumes are not adopted physical geometry.\\
A7 & Any source coframe, internal Lorentzian form, transition representation,
determinant bound, or robustness rule is conditional additional data. &
Current ontology does not derive those objects.\\
A8 & Target geometry, matter, detectors, benchmark success, generated
derivatives, registries, roles, and validators are not scientific premises. &
No target fitting or process-authority smuggling is allowed.\\
\bottomrule
\end{longtable}

\section{Signature and nondegeneracy countermodels}

Work at a source point with an adapted basis
\((e_0,e_1,e_2,e_3)\), oriented ray \(L=[e_0]_+\), standard positive density
\(\mu_0=|e^0\wedge e^1\wedge e^2\wedge e^3|\), and constant scalar phases.

\begin{theorem}[Signature nonselection]
\label{thm:signature}
The symmetric forms
\[
 q_L=\diag(-1,1,1,1),
 \qquad
 q_S=\diag(-1,-1,1,1)
\]
carry identical \(\mathcal D_{\rm cal}\) input and have
\(|\det q_L|=|\det q_S|=1\), but their inertias are respectively
\((1,3)\) and \((2,2)\). Hence the admitted source package does not select
Lorentzian signature, even if determinant magnitude one is imposed as an
additional numerical comparison constraint.
\end{theorem}

\begin{proof}
The ray, annihilator, forward order, independent density, and scalar phases
are defined without either bilinear form and are therefore identical. Direct
diagonalization gives one negative eigenvalue for \(q_L\) and two for
\(q_S\). Their determinants are \(-1\) and \(+1\), respectively. Equal
absolute determinant therefore does not determine inertia.
\end{proof}

\begin{proposition}[Nondegeneracy is not certified]
\label{prop:degeneracy}
For \(\varepsilon>0\),
\[
 q_\varepsilon=\diag(-1,1,1,\varepsilon^2)
\]
is Lorentzian and nondegenerate, but it converges to the degenerate form
\(q_0=\diag(-1,1,1,0)\) as \(\varepsilon\to0\), with
\(\det q_\varepsilon=-\varepsilon^2\). The unchanged source package carries
no lower bound separating this family from the singular boundary.
\end{proposition}

\begin{proof}
The eigenvalues and determinant are displayed explicitly. The independent
density \(\mu_0\) is part of the input for every \(\varepsilon\), but P6-T04
does not identify it with \(|\operatorname{vol}_{q_\varepsilon}|\).
Consequently \(\mu_0\) supplies no metric-determinant certificate. A
conditional law such as \(|\det q|\geq\delta>0\) in a source-derived coframe
would separate the singular boundary, but no such law or coframe is derived.
\end{proof}

\section{A source-naturality obstruction}

A source reconstruction \(F\) is natural only if every automorphism \(A\) of
the complete input satisfies
\[
 F(A\cdot\mathcal D_{\rm cal})=A^*F(\mathcal D_{\rm cal}).
\]
At a fixed input, every input automorphism must therefore preserve the
selected output.

\begin{theorem}[Volume-preserving source naturality obstruction]
\label{thm:naturality}
On \(\mathbb R^4\), take the maximally symmetric datum above and the
determinant-one diagonal source automorphisms
\[
 A(a,b,c)=
 \diag\!\left(e^a,e^b,e^c,e^{-a-b-c}\right).
\]
These maps preserve \(L\), its positive orientation and forward order,
\(\Ann(e_0)\), \(\mu_0\), and the constant phases. The only symmetric
bilinear form satisfying
\[
 A(a,b,c)^TqA(a,b,c)=q
 \quad\hbox{for all }a,b,c\in\mathbb R
\]
is \(q=0\). Therefore no nonzero, and hence no nondegenerate Lorentzian,
bilinear form can be naturally selected on this admitted datum.
\end{theorem}

\begin{proof}
The determinant is
\(e^{a+b+c-a-b-c}=1\), so \(\mu_0\) is preserved. Each map rescales
\(e_0\) by a positive factor and preserves its oriented ray and line order.
The four weights of the representation are
\[
 w_0=(1,0,0),\quad
 w_1=(0,1,0),\quad
 w_2=(0,0,1),\quad
 w_3=(-1,-1,-1).
\]
For a component \(q_{ij}\), invariance gives
\[
 e^{(w_i+w_j)\cdot(a,b,c)}q_{ij}=q_{ij}
\]
for every \((a,b,c)\). No pair \(w_i+w_j\) is zero, including the diagonal
pairs. Thus every \(q_{ij}\) vanishes.
\end{proof}

\begin{corollary}
This theorem is a scoped obstruction for the exact admitted datum and
naturality requirement. It is not a theorem that no richer source ontology
can select geometry. Additional source-derived coframes, constitutive
tensors, discriminators, or transition data can shrink the automorphism group
and reopen the same milestone.
\end{corollary}

\section{Conditional source-coframe gluing theorem}

\begin{theorem}[Conditional Lorentz-coframe gluing]
\label{thm:gluing}
Suppose a source-side construction independently supplies, on each source
domain \(U\), a coframe column
\(\vartheta_U=(\vartheta_U^0,\ldots,\vartheta_U^3)^T\) and a fixed
nondegenerate internal form \(\eta\) of Lorentzian inertia. Define
\[
 q_U=\vartheta_U^T\eta\,\vartheta_U.
\]
If on overlaps
\(\vartheta_V=\Lambda_{VU}\vartheta_U\), then \(q_V=q_U\) exactly when
\[
 \Lambda_{VU}^T\eta\Lambda_{VU}=\eta.
\]
Consistent triple overlaps require
\[
 \Lambda_{WV}\Lambda_{VU}=\Lambda_{WU},
\]
and inverse overlaps require
\(\Lambda_{UV}=\Lambda_{VU}^{-1}\).
\end{theorem}

\begin{proof}
Substitution gives
\[
 q_V=\vartheta_U^T\Lambda_{VU}^T\eta\Lambda_{VU}\vartheta_U.
\]
Because \(\vartheta_U\) is a coframe, equality with \(q_U\) is equivalent to
the displayed Lorentz-preservation identity. Applying two overlap relations
in succession gives
\(\vartheta_W=\Lambda_{WV}\Lambda_{VU}\vartheta_U\); equality with the direct
relation is the cocycle identity. Reversing an overlap gives the inverse
identity.
\end{proof}

The theorem states a genuine positive conditional result. It does not derive
\(\vartheta_U\), \(\eta\), their source locality, a signature discriminator,
the overlap maps, their cocycle, or finite-variation stability. Supplying
these objects would be a proposal-only source-extension law unless separately
derived from current ontology.

\begin{proposition}[Conditional time orientation]
If a connected source region is already equipped with a Lorentzian \(q\) and
the oriented field \(V\) is everywhere \(q\)-timelike, then \([V]_+\)
selects one of the two timelike-cone components and hence a time orientation.
The current input derives neither \(q\) nor \(q(V,V)<0\), so it does not yet
establish a physical time orientation.
\end{proposition}

\section{Precise obstruction and reopening burden}

The exact disposition is
\texttt{OBST-P6T05-SIGNATURE-COVARIANCE-NATURALITY-001}. Identical admitted
source carriers allow different signatures and approach degeneracy, while a
determinant-one automorphism subgroup of one admissible datum forces every
natural symmetric bilinear form to vanish. Current ontology does not derive:
\begin{enumerate}
\item a source coframe or equivalent transverse comparison object;
\item a Lorentzian signature selector or internal form;
\item a metric-volume identity or uniform determinant lower bound;
\item a source-derived Lorentz transition representation with inverse and
cocycle laws; or
\item a finite-variation compatibility and robustness theorem.
\end{enumerate}

A compact proposal-only burden template is therefore:
derive source-local coframes \(\vartheta_U=\mathfrak F_U(\mathcal D)\), a
source discriminator fixing the inertia of \(\eta\), source-induced
\(\Lambda_{VU}\in O(\eta)\) satisfying inverse and cocycle identities, and a
nondegeneracy/robustness certificate. This is \emph{not} a
canonical-ontology candidate or adoption. Any adoption or ontology edit
remains human-gated.

The unchanged route from \(\mathcal D_{\rm cal}\) to signature,
nondegeneracy, and covariant metric selection is locally frozen.
This is not a theorem that the project can never derive those objects. Materially richer
source-side continuation remains open. Current adoption is
\texttt{blocked\_adoption\_open\_continuation}.

\section{Distance-to-GR disposition}

P6-T05 meets its theorem-or-precise-obstruction burden within the exact input
tested. It sharpens Gate B: a source-derived metric would need not merely a
tensorial formula, but a signature discriminator, nondegeneracy certificate,
natural source coframe or equivalent representation, Lorentz-valued
transitions, cocycle and inverse checks, and finite-variation robustness.
Gate B is not passed. The historical scoped \(g_{\rm eff}\) object remains
unchanged. Physical causal geometry, unscoped \(g_{\rm eff}\), matter
coupling, Einstein equations, benchmark promotion, proof authority,
publication, push, and completed derivation remain blocked.

\section{Project source note}

The fixed project basis is the registered P6-T03 and P6-T04 theorem packets,
canonical target-side geometry dictionary, historical scoped source-extension
\(g_{\rm eff}\) candidate and Gate Chair record, v21 implementation plan,
Distance-to-GR burden map, and handoff-0882. These sources are held at the
exact hashes recorded in the accompanying machine-readable specification.
The inertia, determinant, automorphism-weight, and transition calculations are
reproduced directly; no external physical datum is used.

\end{document}
